\phantomsection
\addcontentsline{lot}{section}{\protect\numberline{}Tabel 59 - Jumlah terduga tuberkulosis ,kasus tuberkulosis, kasus tuberkulosis anak dan CNR}
\ra{1.15}

%Table generated by Excel2LaTeX from sheet '56'
{\centering
\begin{tabular}{rY{9em}Y{8em}rZ{4em}Z{4em}Z{4em}rr}
    \multicolumn{9}{l}{Tabel 59}\\
    \multicolumn{9}{c}{JUMLAH TERDUGA TUBERKULOSIS ,KASUS TUBERKULOSIS, KASUS TUBERKULOSIS ANAK,}\\
    \multicolumn{9}{c}{DAN \emph{TREATMENT COVERAGE} (TC) MENURUT JENIS KELAMIN, KECAMATAN, DAN PUSKESMAS}\\
    \multicolumn{9}{c}{KABUPATEN BELITUNG TIMUR}\\
    \multicolumn{9}{c}{TAHUN \tP}\\
    \toprule
    \multirow[t]{3}{*}{NO} & \multirow[t]{3}{*}{KECAMATAN} & \multirow[t]{3}[0]{*}{PUSKESMAS} & \multirow[t]{3}{13em}{\small\raggedleft JUMLAH TERDUGA TUBERKULOSIS YANG MENDAPATKAN PELAYANAN SESUAI STANDAR} & \multicolumn{3}{X{14em}}{JUMLAH SEMUA KASUS TUBERKULOSIS} & \multirow{3}{10em}{\raggedleft JUMLAH KASUS TB SENSITIF OBAT (SO) YANG MEMULAI PENGOBATAN} & \multirow{3}{12em}{\small\raggedleft JUMLAH KONTAK SERUMAH YANG MENDAPATKAN TERAPI PENCEGAHAN TUBERKULOSIS (TPT)}\\
    %\multirow[vpos]{nrows}[bigstruts]{width}[vmove]{text}
    \cmidrule{5-7}
    & & & & L & P & L+P & & \\
    \midrule
    \emph{1} & \emph{2} & \emph{3} & \emph{4} & \emph{5} & \emph{6} & \emph{7} & \emph{8} & \emph{9}\\
    \midrule
    1 & Manggar           & Manggar                 &       &     &       &       &       &       \\
    2 & Damar             & Mengkubang              &       &     &       &       &       &       \\
    3 & Kelapa Kampit     & Kelapa Kampit           &       &     &       &       &       &       \\
    4 & Gantung           & Gantung                 &       &     &       &       &       &       \\
    5 & Simpang Renggiang & Renggiang               &       &     &       &       &       &       \\
    6 & Simpang Pesak     & Simpang Pesak           &       &     &       &       &       &       \\
    7 & Dendang           & Dendang                 &       &     &       &       &       &       \\
    \midrule                                        
    \multicolumn{3}{l}{JUMLAH KAB.}                 & 2.332 & 178 &     98 &   276 &    46  &\\
    \midrule
    \multicolumn{3}{l}{JUMLAH TERDUGA TUBERKULOSIS} & 2.039 & & & & & \\
    \multicolumn{3}{Y{20em}}{\small PERSENTASE ORANG TERDUGA TUBERKULOSIS MENDAPATKAN PELAYANAN TUBERKULOSIS SESUAI STANDAR} & 114,37 & & & & &\\
    \midrule
    \multicolumn{6}{l}{PERKIRAAN INSIDEN TUBERKULOSIS (DALAM ABSOLUT)} & 419 &  & \\
    \multicolumn{6}{l}{CAKUPAN PENEMUAN KASUS TUBERKULOSIS  (\%)} & 419 &  & \\
    \midrule
    \multicolumn{7}{l}{KASUS TUBERKULOSIS SENSITIF OBAT (SO)} & 65,87 & \\
    \multicolumn{7}{l}{PERSENTASE PASIEN TB SO YANG MEMULAI PENGOBATAN (\%)} & 65,87 & \\
    \midrule
    \multicolumn{8}{l}{PERKIRAAN JUMLAH KONTAK SERUMAH YANG DIBERIKAN TERAPI PENCEGAHAN TUBERKULOSIS (TPT)} & 91,49 \\
    \multicolumn{8}{l}{CAKUPAN PEMBERIAN TERAPI PENCEGAHAN TB PADA KONTAK SERUMAH} & 91,49 \\
    \bottomrule
\end{tabular}%

}

\vfill
Sumber: Subkoordinator Pengendalian Penyakit Menular\par
